\chapter{Education Pipelines}\label{ch:in:education-pipelines}

In the year to June 2024 US universities awarded 4\,584 master's degrees coded as financial mathematics and 870 coded as
financial analytics, a code that did not exist before 2020. In 2010 the first code counted 305 master's degrees. The
same year's completions survey counts 25\,543 master's degrees in computer science and 1\,760 research doctorates in
physics, and the industry hires from all of them. This chapter counts the degrees that feed the industry, reads two
specialised programmes' \omterm{def:in:education-pipelines:placement}{placement reports} for what they say and do not say, and looks at doctorates, competitions and
learning outside a programme.

\section{Degrees that feed each kind of firm}

\begin{definition}[Classification of Instructional Programs]\label{def:in:education-pipelines:cip}
The \emph{Classification of Instructional Programs}\index{Classification of Instructional Programs} (CIP) is the US
taxonomy of fields of study, with six-digit codes and definitions, revised every ten years; institutions report their
awards to the federal completions survey by CIP code and award level.
\end{definition}

The federal completions survey counts awards, not people or curricula: each programme chooses its code. The code for
financial mathematics, unchanged since the 2010 classification, is defined as ``the application of mathematics and
statistics to the finance industry, including the development, critique, and use of various financial models''; the
2020 revision added financial analytics, ``financial big data modeling''. Neither code covers every financial engineering
programme, and some of the growth in the first may reflect programmes adopting it; the counts are therefore a floor on
specialised financial degrees, and a guide to the trend.

\begin{dated}{2024-06}{US degrees in the industry's fields}\label{dat:in:education-pipelines:degrees}
Awards to first majors in the year to June 2024 (IPEDS): financial mathematics 611 bachelor's, 4\,584 master's, 24
research doctorates; financial analytics 553, 870, 4; statistics 3\,379, 2\,716, 454; mathematics 14\,333, 2\,230,
1\,165; computer science 44\,627, 25\,543, 1\,628; physics 6\,008, 1\,934, 1\,760. Master's degrees in financial mathematics:
305 in 2010, 773 in 2014, 3\,554 in 2019.
\end{dated}

\begin{omfigure}
\begin{tikzpicture}
\begin{axis}[width=10.5cm,height=5.8cm,xmin=2009,xmax=2025,ymode=log,ymin=100,ymax=40000,
  xtick={2010,2014,2019,2024},xticklabel style={/pgf/number format/1000 sep=},xlabel={\small year of award (ending June)},
  ylabel={\small master's degrees (log scale)},tick label style={font=\footnotesize},grid=major,grid style={gray!20},
  table/col sep=comma,unbounded coords=jump,log basis y=10,
  legend style={font=\scriptsize,at={(0.5,-0.3)},anchor=north,legend columns=3,legend cell align=left,
  /tikz/every even column/.append style={column sep=6pt}}]
\addplot[omThm,thick,mark=square*] table[x=year,y=financial_mathematics]{figdata/industry/29-education-pipelines/masters.csv};
\addplot[omThm!50,thick,mark=triangle*] table[x=year,y=financial_analytics]{figdata/industry/29-education-pipelines/masters.csv};
\addplot[omDef,thick,mark=*] table[x=year,y=statistics]{figdata/industry/29-education-pipelines/masters.csv};
\addplot[black!70,thick,mark=o,mark options={solid}] table[x=year,y=mathematics]{figdata/industry/29-education-pipelines/masters.csv};
\addplot[omDef!50,thick,dashed,mark=diamond*,mark options={solid}] table[x=year,y=computer_science]{figdata/industry/29-education-pipelines/masters.csv};
\addplot[black!40,thick,mark=x] table[x=year,y=physics]{figdata/industry/29-education-pipelines/masters.csv};
\legend{financial mathematics,financial analytics,statistics,mathematics,computer science,physics}
\end{axis}
\end{tikzpicture}

\omcaption{US master's degrees awarded to first majors in six fields, years ending June 2010, 2014, 2019 and 2024
(IPEDS completions). Financial analytics is a 2020 code and appears once. Data:
\texttt{data/industry/ipeds\_completions.csv}, through \texttt{in\_edu.completions}.}\label{fig:in:education-pipelines:masters}
\end{omfigure}

Different firms draw on different degrees. Research roles (chapter~17) draw on doctorates and master's degrees in
mathematics, statistics, physics and computer science; engineering roles (chapters~19 and~20) on computer science at every
level; bank quant roles (chapter~18) and structuring (chapter~23) on the specialised financial master's degrees as well as
on doctorates. The specialised degrees are small beside the general ones (\cref{fig:in:education-pipelines:masters}):
financial mathematics' master's degrees were less than a fifth of computer science's in 2024.

\section{The specialised master's programmes: size, cost and placement}

\begin{definition}[Placement report]\label{def:in:education-pipelines:placement}
A \emph{placement report}\index{placement report} is a degree programme's own account of what its graduates did after the
degree: how many sought work, how many accepted offers by a given date, where, and the pay reported by those who reported
it, under a reporting standard the programme chooses or none.
\end{definition}

A \omterm{def:in:education-pipelines:placement}{placement report} is evidence from an interested party, and one programme says so itself: ``There is no common standard
for reporting placement information among professional financial engineering master's programs, and the reports of such
programs are not audited.'' The reader's questions are the denominators: of how many students, how many sought work, how
many accepted offers, and how many reported pay.

\begin{dated}{2026-01}{Two placement reports}\label{dat:in:education-pipelines:placement}
UC Berkeley Master of Financial Engineering, class of 2025: 70 students looking for jobs, 67 accepted offers; median \omterm{def:in:how-pay-works:base}{base
salary} \$150\,000; the report does not say how many reported pay. Carnegie Mellon MS in Computational Finance, class of
2026 (graduated December 2025): 99 students, 98 seeking work, 98 with accepted offers within three months, counting
short-term employment under its reporting standard; median \omterm{def:in:how-pay-works:base}{base salary} \$145\,000 from 89 students reporting.
\end{dated}

\begin{method}[What a median from part of a cohort can say]\label{met:in:education-pipelines:bounds}
If $r$ of $n$ graduates report pay, the median of all $n$ lies between the reported distribution's quantiles
$\max\bigl(0,(n/2-(n-r))/r\bigr)$ and $\min\bigl(1,(n/2)/r\bigr)$, whatever the $n-r$ non-reporters earn: all below the
reporters at one extreme, all above at the other.
\end{method}

For the second report, 89 of 98 reported: the median of all 98 lies between the reported pay's 44.9th and 55.1st
percentiles, a narrow range. For the first, the number reporting is not stated, and no such bound is possible. Neither
report says what non-reporters earned or why they did not report, and a rate of 96\% or 100\% ``placed'' depends on what
counts as a job and by when. Cost and length of study belong in the same comparison.

\section{Doctorates and the academic route}

The US awarded 1\,760 research doctorates in physics, 1\,628 in computer science, 1\,165 in mathematics and 454 in
statistics in the year to June 2024 (\cref{fig:in:education-pipelines:levels}), together about as many as the master's
degrees in financial mathematics. Doctorates enter research roles at a level above new graduates (chapter~28), often
after an \omterm{def:in:career-paths-and-moves:entry}{internship} during the degree. The academic route runs in the other direction too: researchers leave the
industry for faculty positions, and faculty consult for or join firms; neither flow is counted by any public source this
book found.

\begin{omfigure}
\begin{tikzpicture}
\begin{axis}[xbar,area legend,width=10.5cm,height=5.8cm,xmode=log,xmin=1,xmax=100000,log origin=infty,ymin=0.4,ymax=6.6,
  bar width=5pt,ytick=data,yticklabels from table={figdata/industry/29-education-pipelines/levels.csv}{field},
  yticklabel style={font=\footnotesize},xticklabel style={font=\footnotesize},
  xlabel={\small awards, year to June 2024 (log scale)},grid=major,grid style={gray!20},table/col sep=comma,
  legend style={legend cell align=left,font=\scriptsize,at={(0.5,-0.36)},anchor=north,legend columns=3,/tikz/every even column/.append style={column sep=6pt}}]
\addplot[fill=black!25,draw=none] table[x=bachelor,y=pos]{figdata/industry/29-education-pipelines/levels.csv};
\addplot[fill=omDef,draw=none] table[x=master,y=pos]{figdata/industry/29-education-pipelines/levels.csv};
\addplot[fill=omThm,draw=none] table[x=doctorate,y=pos]{figdata/industry/29-education-pipelines/levels.csv};
\legend{bachelor's,master's,research doctorate}
\end{axis}
\end{tikzpicture}

\omcaption{US awards in six fields by level, year to June 2024 (IPEDS, first majors). Data:
\texttt{data/industry/ipeds\_completions.csv}, through \texttt{in\_edu.completions}.}\label{fig:in:education-pipelines:levels}
\end{omfigure}

\section{Competitions and what they signal}

Firms that hire quantitative graduates read competition results as evidence of problem-solving under time pressure. The
best known is the International Mathematical Olympiad, for students ``Under 20 years old and not enrolled in
post-secondary education'': national teams of up to six, ``2 days, 3 problems per day, 4.5 hours each day'', each problem
worth 7 points. University programming and mathematics competitions play a similar role later. What a result signals is
narrow: speed and depth on well-posed problems. Research asks for problems no one has posed, data that lie, and patience
over months; Book~18 discusses how interviews try to measure both.

\section{Learning outside a programme}

Much of what the industry uses is learnt outside a degree: programming in production, statistics on real data, the
markets themselves. The routes are open: textbooks and open courses, open-source libraries (chapter~20's public
engineering), data competitions, and the running project of this series, whose code accompanies every book. A candidate
without a specialised degree shows the same evidence a programme would: work that can be read and rerun.

\section{Tutorial: how many quants does the world train?}\label{tut:in:education-pipelines:tut}

\begin{tutorial}
\textbf{Goal.} Count the US degrees in the industry's fields, read two \omterm{def:in:education-pipelines:placement}{placement reports} for their response, and put
the specialised degrees beside the filings of chapter~14.
\textbf{End state:} \cref{fig:in:education-pipelines:masters,fig:in:education-pipelines:levels} and the placement bounds.
\begin{tutsteps}
\item \textbf{Completions.} \texttt{firm.edupipe.read\_ipeds} streams one year's zip and totals first-major awards by code
and level (\cref{lst:in:education-pipelines:ipeds}); \texttt{in\_ipeds\_derive.py} runs it for four years.
\omcode{code/firm/edupipe/firm_edupipe.py}{23}{36}{Awards by programme code and level from one year's completions file.}\label{lst:in:education-pipelines:ipeds}
\item \textbf{\omterm{def:in:education-pipelines:placement}{Placement reports}.} \texttt{PlacementReport} records the denominators; \texttt{median\_quantile\_bounds}
gives the bound of the method (\cref{lst:in:education-pipelines:bounds}).
\omcode{code/firm/edupipe/firm_edupipe.py}{59}{66}{The quantiles between which a cohort's median must lie.}\label{lst:in:education-pipelines:bounds}
\item \textbf{The ratio.} Master's degrees in the two financial codes against the fiscal 2025 filings in the
\omterm{def:in:quant-researcher:def}{quantitative researcher} and quant developer families.
\end{tutsteps}
\end{tutorial}

The result: 5\,454 master's degrees in 2024 against 1\,466 filings in fiscal 2025, a ratio of 3.72. The two numbers
count different things (degrees of all nationalities against visa filings by a sample of employers) and the ratio is a
scale, not a placement rate.

\textbf{What to change next.} Add the other years' files and the UK's counts where openly licensed; split completions
by citizenship (the survey reports nonresident students); collect more \omterm{def:in:education-pipelines:placement}{placement reports} with their denominators.

\section{Build: firm.edupipe}\label{bld:in:education-pipelines:edupipe}

\begin{build}
\bfield{Purpose} Count degrees by field and level from the federal survey, and record \omterm{def:in:education-pipelines:placement}{placement reports} with their
denominators.
\bfield{Interface} \texttt{firm.edupipe}: \texttt{LEVELS}; \texttt{read\_ipeds(zip\_path, cips, levels)};
\texttt{PlacementReport(programme, cohort, students, seeking, accepted, reporting, median\_base, standard, note)} with
\texttt{placement\_rate}, \texttt{reporting\_rate} and \texttt{median\_quantile\_bounds()}.
\bfield{Rules} First majors only; the file is streamed; a report that does not state its reporting count gives no bound.
\bfield{Acceptance tests} \texttt{code/firm/edupipe/tests/}: second majors and other codes are excluded and zero-padded
levels are read; 89 of 98 reporting bound the median between the 44.9th and 55.1st reported percentiles.
\bfield{Stretch} HESA counts; the nonresident share of completions; a placement-report parser.
\end{build}

\begin{omsources}
\item National Center for Education Statistics, IPEDS completions and the Classification of Instructional Programs 2020.
\item UC Berkeley MFE and Carnegie Mellon MSCF employment reports.
\item International Mathematical Olympiad, general information.
\item Chapter~14's tables from the Department of Labor's LCA files.
\end{omsources}

\section{Exercises}

\begin{exercise}[$\star$]\label{exo:in:education-pipelines:1}
By what factor did master's degrees in financial mathematics grow between 2010 and 2024?
\end{exercise}

\begin{exercise}[$\star$]\label{exo:in:education-pipelines:2}
What share of the Berkeley class looking for jobs accepted an offer?
\end{exercise}

\begin{exercise}[$\star$]\label{exo:in:education-pipelines:3}
How many points can an IMO contestant score at most?
\end{exercise}

\begin{exercise}[$\star\star$]\label{exo:in:education-pipelines:4}
If 60 of 98 had reported pay in the second report, between which reported percentiles would the cohort's median lie?
\end{exercise}

\begin{exercise}[$\star\star$]\label{exo:in:education-pipelines:5}
Why is the count of financial mathematics degrees a floor on specialised financial degrees, and how could it overstate
the trend?
\end{exercise}

\begin{exercise}[$\star\star$]\label{exo:in:education-pipelines:6}
Compare the research doctorates in the four general fields with the financial master's degrees in 2024.
\end{exercise}

\begin{exercise}[$\star\star\star$]\label{exo:in:education-pipelines:7}
\emph{Coding.} Using \texttt{read\_ipeds}, compute the growth of computer science master's degrees between 2019 and 2024,
and compare it with financial mathematics'.
\end{exercise}

\begin{exercise}[$\star\star\star$]\label{exo:in:education-pipelines:8}
\emph{Find the flaw.} ``Programme A reports a median of \$150\,000 and programme B \$145\,000, so A places its graduates
better.''
\end{exercise}

\section{Problem: How Many Quants Does the World Train?}

\begin{problem}[{Weekend problem --- how many quants does the world train?}]\label{pb:in:education-pipelines:1}
A university considers launching a financial engineering master's and asks how large the market for its graduates is.

\medskip\noindent\textbf{Part I --- The codes.}
\begin{enumerate}
\item Define the \omterm{def:in:education-pipelines:cip}{Classification of Instructional Programs}.
\item Which codes cover financial mathematics and financial analytics, and how are they defined?
\item What does the completions survey count, and over which period?
\item Why do first majors matter?
\item Why is the count a floor?
\end{enumerate}

\medskip\noindent\textbf{Part II --- The numbers.}
\begin{enumerate}[resume]
\item Give the master's degrees in financial mathematics for the four years.
\item Give the 2024 awards by level in the six fields.
\item How large are the specialised degrees beside computer science?
\item Give the research doctorates in the four general fields.
\item Which roles draw on which degrees?
\end{enumerate}

\medskip\noindent\textbf{Part III --- Placement.}
\begin{enumerate}[resume]
\item Define a \omterm{def:in:education-pipelines:placement}{placement report}.
\item Give the two reports' denominators and medians.
\item State the bound on a cohort's median and apply it.
\item What does neither report say?
\item What does a competition result signal, and what not?
\end{enumerate}

\medskip\noindent\textbf{Part IV --- The verdict.}
\begin{enumerate}[resume]
\item State the \emph{named result}: annual US master's completions in the financial programme codes, and their ratio to
the yearly count of LCA filings in quantitative titles at finance employers.
\item Why is the ratio not a placement rate?
\item What would the university's graduates compete with?
\item What should its \omterm{def:in:education-pipelines:placement}{placement report} publish?
\item In two sentences, answer the university.
\end{enumerate}
\end{problem}

\section{Interview questions}

\begin{interviewq}[$\star$ \iqroles{researcher}]\label{iq:in:education-pipelines:1}
What did your thesis teach you that a course could not?
\end{interviewq}

\begin{interviewq}[$\star$ \iqroles{researcher, developer}]\label{iq:in:education-pipelines:2}
Show us a piece of work outside your coursework and explain one decision in it.
\end{interviewq}

\begin{interviewq}[$\star\star$ \iqroles{researcher}]\label{iq:in:education-pipelines:3}
A report says 100\% of a class was placed. What would you ask before believing it?
\end{interviewq}

\begin{interviewq}[$\star\star$ \iqroles{researcher}]\label{iq:in:education-pipelines:4}
How would you estimate the number of people who become \omterm{def:in:quant-researcher:def}{quantitative researchers} each year?
\end{interviewq}

\begin{interviewq}[$\star\star$ \iqroles{developer}]\label{iq:in:education-pipelines:5}
You taught yourself C++. How would you show us you can write production code?
\end{interviewq}

\begin{interviewq}[$\star\star\star$ \iqroles{researcher}]\label{iq:in:education-pipelines:6}
A counted series jumps when a new classification code is introduced. How do you tell a real change from a recoding?
\end{interviewq}
